%% ma: L11-B16-0058
%% lop: 11
%% chuong: 5
%% bai: 16
%% dang: 11.16.4
%% loai: DS
%% muc: [NB, TH, VD, VD]
%% dap_an: "SĐSĐ"
%% nguon: "tai-lieu-hoc-tap-mon-toan-lop-11-theo-chuong-trinh-giao-duc-pho-thong-2018.pdf (Lớp Toán Thầy Duy, Phạm Lê Duy, THPT Chu Văn An), Chương 3 Bài 2, Phần 2, Câu 2 (tr.375)"
%% kien_thuc: "Lớp 11 Bài 16 (giới hạn của hàm số)"
%% trang_thai: cho-duyet
%% ly_do: "Dạng mới đề xuất (11.16.x), chờ giáo viên duyệt dạng; tài liệu gốc không có lời giải mẫu, đáp án do kho tự giải và kiểm bằng sympy."
\begin{ex}
Cho $\lim\limits_{x\to-\infty}\dfrac{a\sqrt{x^2+1}+2023}{x+2024}=\dfrac12$; $\lim\limits_{x\to+\infty}\left(\sqrt{x^2+bx+1}-x\right)=2$. Khi đó
\choiceTF
{$a>0$.}
{\True $b>0$.}
{$a>b$.}
{\True $P=4a+b=2$.}
\loigiai{\sloppy Với $x<0$: $\sqrt{x^2+1}=-x\sqrt{1+\frac1{x^2}}$, nên $\dfrac{a\sqrt{x^2+1}+2023}{x+2024}\to-a=\dfrac12\Rightarrow a=-\dfrac12$. Và $\sqrt{x^2+bx+1}-x=\dfrac{bx+1}{\sqrt{x^2+bx+1}+x}\to\dfrac b2=2\Rightarrow b=4$.
a) $a=-\dfrac12<0$. Sai.
b) $b=4>0$. Đúng.
c) $a=-\dfrac12<4=b$. Sai.
d) $4a+b=-2+4=2$. Đúng.}
\end{ex}
